%
% Revision 30 Nov 2023
%
%%%%%%%%%%%%%%%%%%%%%%%%%%%%%%%%%%%%%%%%%%%%%%%%%%%%%%%%%%%%%%%%%%%%%%%%%%%%%%%%
%2345678901234567890123456789012345678901234567890123456789012345678901234567890
%        1         2         3         4         5         6         7         8
% DOCUMENT CLASS
\documentclass[a4paper,12pt]{Classes/RoboticsLaTeX}

\usepackage{amsmath}
\usepackage{amsfonts}
\usepackage{algorithm}
\usepackage{algorithmic}
\usepackage{multirow}
\usepackage{colortbl}
\usepackage{color}
\usepackage[table]{xcolor}
\usepackage{epigraph}
\usepackage{graphicx}
%\usepackage{subfigure}
\usepackage{caption}
\usepackage{subcaption}
\usepackage{hyperref}
\usepackage{tabularx}
\usepackage{float}
\usepackage{longtable}
\usepackage[pdftex]{graphicx}
\usepackage{pdfpages}
\usepackage{pdflscape}
\usepackage[acronym,toc]{glossaries}
\usepackage{setspace}
\usepackage[utf8]{inputenc}
\usepackage[table]{xcolor}

\setstretch{1.5}
% correct bad hyphenation
\hyphenation{op-tical net-works semi-conduc-tor}
\hyphenation{par-ti-cu-lar mo-du-le ge-stu-re}

\title{\Large{Assignment Predictions for Machine Reassignment Problem}}
\author{Piotr Bielski}
\collegeordept{School of Computer Science}
\university{University of Galway}
\crest{
\includegraphics[width=135mm]{Figures/University_Of_Galway_Logo__Positive_Landscape.png}}
\supervisor{Takfarinas Saber}
\degree{MSc in Computer Science (Artificial Intelligence - Online)}
\degreedate{2025}

%%%%%%%%%%%%%%%%%%%%%%%%%%%%%%%%%%%%%%%%%%%%%%%%%%%%%%%%%%%%%%%%%%%%%%%%%%%%%%%%
%%% uncomment if glossary needed, see examples in file
%\makeglossaries
%\loadglsentries{glossary}

\begin{document}
    
\begin{spacing}{1}
    \maketitle
\end{spacing}
	
% add an empty page after title page
\newpage\null\thispagestyle{empty}\newpage
	
% set the number of sectioning levels that get number and appear in the contents
\setcounter{secnumdepth}{3}
\setcounter{tocdepth}{3}
	
\frontmatter
	
\textbf{DECLARATION} 
I, Piotr Bielski, hereby declare that this thesis, titled ``Solutions Analysis to Machine Reassingnment Problem'', and the work presented in it are entirely my own except where explicitly stated otherwise in the text, and that this work has not been previously submitted, in part or whole, to any university or institution for any degree, diploma, or other qualification. 
\newline
    
\begin{tabular}{@{}p{.5in}p{4in}@{}}
    Signature: & ~~\hrulefill \\
\end{tabular}
\newpage
	
%%%% uncomment if acknowledgements needed
%\textbf{Acknowledgement}
%
%
%\newpage\textbf{}



\begin{abstracts}

The reliance on cloud and high-performance computing platforms is becoming increasingly more widespread due to the rising demand for more computational resources across wide range of applications.
However, the supporting infrastructure requires significant electrical power and negatively impacts the environment. Additionally, at such large scales, any inefficiencies in managing compute resources results in unnecessary and significant costs incurred by the compute providers or consumers.
Therefore, to reduce environmental impact while at the same time reducing operational costs, the implementation of efficient resource management systems is essential. 
This report aims to investigate challenges in efficiently managing compute resources, research into optimization algorithms, specifically solutions for solving the ROADEF/EURO 2012 Challenge (Machine Reassignment Problem), common techniques and reassignment patterns among best performing solutions, and synthesis of characteristics for the future development of optimal reassignment algorithms.
\newline
\newline
\textbf{Keywords:} Machine Reassignment, Resource Management, Optimization, ROADEF, MRP

\end{abstracts}

\tableofcontents
\listoffigures
\listoftables
\printglossary[title=List of Acronyms,type=\acronymtype]
\mainmatter



\chapter{Introduction}\label{chap:introduction}

\textbf{(What) Topic:} High-performance computing and cloud platforms provide wide range of services in the form of Infrastructure-as-a-Service (IaaS), Platform-as-a-Service (PaaS), Data-as-a-Service (DaaS), which provide easy access to additional storage and compute resources. These services are in growing demand due to technological advancements in the recent years, resulting in
explosion of resource-intensive scientific and industrial applications
such as; artificial intelligence, SaaS enterprises, adoption of
distributed systems across manufacturing and energy industries, as well
as autonomous systems in the automobile industry - and all that with
data centers as the backbone running this infrastructure. 
\cite{dillonCloudComputingIssues2010}
The significant growth in demand for computational resources is pushing the construction on new large-scale data centers, with an estimation of
Google having around one million servers. 
\cite{bharanySystematicSurveyEnergyEfficient2022}
While the benefits are plentiful across many industries, the data centers are requiring an
increasingly larger supply of electrical power, with an estimation of
already consuming around 1\% of the global energy supply, while still
being on the rise, 
\cite{masanetRecalibratingGlobalData2020}
and having a growing impact on the environment due to CO2 emissions. 
\cite{beloglazovEnergyEfficientResource2010}

\textbf{(Why) Motivation:} Resource management strategies are essential
for efficient utilization of resources, which in turn reduce wastage of
electrical power and operational costs. However, cloud providers are
facing significant difficulties in keeping up with the rate of growing
demand to maintain efficiency and quality-of-service (QoS) standards due
to the complexities of the dynamic nature of the environment, where the
workload (e.g.~network traffic) is not constant but instead fluctuates,
requiring both intelligent and adaptive systems to optimally distribute
workload among the available resources. \cite{dillonCloudComputingIssues2010}
Considerable strategies have already been explored to improve the resource management
efficiency, however, there is still much potential for further
optimizations, and even small improvements can yield significant results
at such scale. \cite{masanetRecalibratingGlobalData2020}

\textbf{(How) Purpose:} In this context, Google presented a complex and
large-scale machine reassignment problem (MRP) for the ROADEF/EURO
contest in 2012 (https://www.roadef.org/challenge/2012/en/index.php)
with the purpose of driving more efforts into research and innovation of
novel optimization algorithms for tackling the major challenges of
efficiently managing resources in dynamic environments.

The goal of this report is to analyze the behavior of multiple different
best performing state-of-the-art optimization algorithms submitted for
solving the ROADEF12 challenge by tracking the reassignment of processes
during the lifecycle of the optimization in an attempt to provide new
knowledge into what types of processes are reassigned by the different
algorithms and to investigate whether there are any patterns in the
reassignments or other characteristics that could be used as guideline
for development of better resource management optimization algorithms.


\newpage
\section{Research Questions}

\begin{enumerate}
\def\labelenumi{\arabic{enumi}.}
\item RQ: What are the state-of-the-art algorithms for tackling the machine reassignment problem and do they have any common characteristics?

\item RQ: What processes are reassigned by the top performing algorithms over the duration of the optimization process and are there any patterns in reassignments?

\item RQ: Can insights into the reassignments be used to create a framework for guiding the design of heuristics towards better selection of processes for reassignments? i.e.~to leverage these patterns to create a set of heuristics/ rules to guide the reassignment process in a more efficient manner.

\item RQ: Can the insights be used in the development of future machine reassignment optimization algorithms?
\end{enumerate}



\newpage
\section{Report Structure}\label{report-structure}

\textbf{Chapter 1: Introduction} lays out the foundation for the
project, presenting the impact of growing reliance on compute resources,
the challenges in efficiently managing these resources, and approaches
in tacking the issue.

\textbf{Chapter 2: Background} provides background information on the
ROADEF/EURO 2012 Challenge (Machine Reassignment Problem), the problem
definition, and common algorithms used for tackling the optimization
problem, including information on reassignment and search strategies.

\textbf{Chapter 3: Related Works} discusses the research objectives in
more detail, the research methods applied when searching for related
works, and an analysis of proposed solutions.

\textbf{Chapter 4: Methodology} presents the research design and
approach, specifically the required metrics for analyzing the evolution
of solution over time, and the design of the tracking system to generate
the dataset from reassignments.

\textbf{Chapter 5: Experiments} contains all necessary steps to
reproduce the experimental results.

\textbf{Chapter 6: Results} investigates the recorded reassignments,
evaluates reassignments among solutions, and identifies reassignment
patterns.



\chapter{Background} \label{chap:background}

The following sections provide background on the general topic and challenges in machine reassignment, an introduction to ROADEF/EURO annual contests, with expanded information on the specific machine reassignment problem contest from 2012, and a detailed problem definition.



\section{Machine Reassignment Problem}

\textbf{Machine Reassignment Problem (MRP)} is a complex combinatorial
optimization problem that involves the distribution of workload
(processes) across a set of machines in a way to optimize the
utilization of available resources on the machines, and is particularly
relevant in dynamic environments, where workloads fluctuate and
intelligent algorithms are needed to improve efficiency of the system.

However, the major challenge stems from the complexity of the problem -
which requires balancing many variables, of which values may change, as
is the case in cloud computing where workloads are continuously received
or completed resulting in the conditions always changing, and such, not
only it is insufficient to run the optimization algorithm once to
distribute the workload, but instead it is an on-going iterative process
which involves constant evaluation of the system and potentially cycles
of steps of reorganizing the processes across machines based on current
workloads and available resources, but also given the significant number
of variables necessary for consideration, e.g.~resources required by the
process (RAM, CPU, DISK, etc.), available resources on each machine, the
total capacity of each machine, or balancing load across machines - the
MRP problem is defined as NP-hard. This classification implies that
finding an optimal solution is computationally not feasible, especially
for instances of realistic size and complexity, and hence often requires
the use of heuristic and metaheuristic approaches to tackle these types
of problems. Such approaches aim to find high-quality solutions in a
reasonable timeframe, but with the trade-off of no guarantee for
obtaining global optimality. \cite{wreford2024}



\section{ROADEF/EURO 2012 Challenge - Google Machine Reassignment Problem}

The ROADEF/EURO contest is jointly organized by French Operational Research and Decision Aid Society (ROADEF) and Association of European Operational Research Society (EURO) and appears regularly since 1999. 
\footnote[1]{https://www.roadef.org/challenge/2022/en/}

The challenges are open to everyone and aim to promote innovation and
collaboration between researchers. Generally, participants are split
into two categories; (1) Junior category for young researchers, and (2)
Senior category for experienced and qualified researchers. The
challenges provide an opportunity to tackle complex industrial
optimization problems and have been received as widely successful.
\cite{murat2016}

As for the specific challenge related to this report, Google presented a
formulation of the machine reassignment problem for the ROADEF/EURO 2012
contest, called the Google Machine Reassignment Problem (``GMRP''),
\footnote{https://roadef.org/challenge/2012/en/} with a set of
benchmarking instances to the research community. The problem is
commonly encountered in dynamic environments such as cloud computing,
high-performance computing platforms, data centers, and other similar
large compute platforms, where workloads can fluctuate and it is
beneficial to reallocate available resources to where they are needed
based on current demand.

Once published, 82 teams registered for the GMRP contest. The first
stage involved the participants to submit their initial solutions to be
evaluated against the smallest dataset, ``A'', of which, only 30 teams
could qualify into the next round. In the second round, the qualified
teams then moved to more realistic set of instances made available to
them, called dataset ``B'', to allow for fine-tuning their algorithms
and prepare for having their solutions run against the final, largest,
and unknown dataset ``X''.

\emph{Could modify above to mention the setup across all dataset
instances was; given initial assignment of processes to machines, the
algorithms run for period of 5 minutes, and upon completion, the output
consisting of final process assignments would be passed through solution
checker to compute to final solution cost, and then evaluated against
results of other participants.}

Overall, the GMRP contest created a competitive environment that
resulted in a diverse array of novel state-of-the-art approaches and
highlights the problem's profound relevance in the industry. The results
and algorithms from the contest offer rich dataset for the analysis of
efficacy of different algorithmic strategies and other explored
domain-specific heuristics that pushed performance of the solutions even
further beyond their original capabilities.



\section{Problem Definition}

GMRP \footnote[2]{https://roadef.org/challenge/2012/en/} is formulated
through three primary components; decision variables, constraints, and
objectives, each consisting for interacting entities. The fundamental
goal of the problem is; given an initial assignment of processes to
machines, to reassign the processes to optimize for the objective
function while adhering to multitude of constraints.

\subsection{Decision Variables}\label{decision-variables}

The fundamental decision variables defined for the problem are;
machines, processes, and the initial assignment, where: \emph{(Could
improve this explanation)}

\begin{itemize}
\item 
    \textbf{Machines} refers to set of available machines, \(\mathcal{M}\), used for the execution of processes, where each machine, \(\textit{m} \in \mathcal{M}\), has a set of pre-defined finite resources such as CPU, RAM, DISK.
\item
    \textbf{Processes} refers to set of processes, \(\mathcal{P}\), that are the fundamental units of work. Each process, \(\textit{p} \in \mathcal{P}\), has specific resource requirements (CPU, RAM, DISK) and consumes these resources off the host when assigned to a machine, \(\mathcal{M}(\textit{p}) = \textit{m}\).
\item
    \textbf{Initial Assignment} defines the initial mapping of processes to machines in a way that may not be optimal, but to satisfy constraint requirements. The initial assignments are denoted as \(\mathcal{M}_0(\textit{p})\).
\end{itemize}



\subsection{Constraints}\label{constraints}

Constraints define a set of limitations to which a solution must adhere and reflect physical restrictions such as available compute resources, process conflicts, dependencies, or the migration of processes over distributed network. Through these boundaries, not only do the constraints define the feasible solution space but also act as an implicit driver for the reassignment process:

\begin{itemize}
\item
  \textbf{Capacity Constraints}: A process cannot be assigned to a
  machine if the machine does not have enough available resources,
  \(\mathcal{R}\). The machine must have enough capacity of all
  resources (CPU, RAM, DISK) to host the process. The capacity of
  resource is defined as \(\mathcal{C} (\textit{m}, \textit{r})\).
\item
  \textbf{Conflict Constraints}: Defines that processes belonging to
  same service must run on distinct machines to avoid conflicts, and
  does not allow processes of the same service to run on the same
  machine, where a service, \(\textit{s} \in \mathcal{S}\), consists of
  processes.
\item
  \textbf{Spread Constraints}: Defines the minimum number of locations,
  \emph{minSpread}, over which processes of a service must be
  distributed, where \emph{location}, \(\textit{l} \in \mathcal{L}\),
  consists of a set of machines.
\item
  \textbf{Dependency Constraints}: Services can have a dependency on
  another service. In such cases, processes of the dependent service
  must run in the neighbourhood of the other service processes. Here, a
  neighbourhood is defined as a set of machines as well,
  \(\textit{n} \in \mathcal{N}\).
\item
  \textbf{Transient Usage Constraints}: While a process is migrated from
  one machine, \textit{m}, to another, \textit{m}', the resources
  required by the process are temporarily consumed on both machines.
  Eventually, once the migration is complete, the resources are freed
  from the initial machine.
\end{itemize}

\subsection{Objectives}\label{objectives}

Objectives are the variables used by algorithms for guiding the
optimization process and defines incurred costs from certain actions,
but also penalties for violating certain conditions:

\begin{itemize}
\item
  \textbf{Load Cost} is incurred for over-utilizing machine resources
  beyond predefined safety capacity threshold, and serves as penalty for
  pushing machine usage beyond preferred operational load, even if still
  within the absolute physical limits. Each machine defines a safety
  capacity per resource, \(\mathcal{SC} (\textit{m}, \textit{r})\), over
  which the penalty is incurred.
\item
  \textbf{Balance Cost} defines an objective to encourage maintaining a
  balance between available resources across machines and incurs a
  penalty for imbalances.
\item
  \textbf{Process Move Cost} accounts for the overhead of moving a
  process from one machine to another and incurs a cost for each process
  reassignment.
\item
  \textbf{Service Move Cost}: Designated to balance the number of moved
  processes across different services. It is defined as maximum number
  of moved processes for any single service.
\item
  \textbf{Machine Move Cost}: Cost of migrating any process to another
  machines. The cost varies depending on specific source and destination
  machines.
\item
  \textbf{Total Objective Cost}: Weighted sum of all individual costs.
\end{itemize}

In summary, the \textbf{Total Objective Cost} is an aggregation of
individual cost components, and effectively transforms the problem from
multi-objective into single-objective optimization problem.



\section{Solution Components: Approaches, Exploration, and Search Mechanisms}

At its core, reassignment is the fundamental action of changing the
allocation of a process from one machine to another, and is also the
only action needed to improve the efficiency of the larger system
through optimization towards better utilization of available resources,
provided the reassignments are controlled using an intelligent
optimization algorithm. Given it is not feasible to run the optimization
process for the entire problem at once, a common approach involves
reducing the scope by considering only a subset of processes to optimize
through reassignments, and iteratively repeating this procedure each
time selecting another subset of processes. For such algorithm, the two
primary components are; subproblem selection and search mechanisms -
these define how the algorithm determines which subset of processes is
selected in each iteration and how it explores the solution space
through possible reassignments for the current subset of processes. This
approach is the basis for two fundamental search paradigms used to
tackle the MRP; local search and large neighborhood search.



\subsection{Levels of Abstraction in Optimization Approaches}

Algorithms are generally categorized based on their approaches \cite{wreford2024} and include:

\begin{itemize}
\item
  \textbf{Heuristic} approaches are simple, problem-specific rules
  designed to quickly find ``good enough'' solutions and rely on domain
  knowledge of the specific problem, e.g.~in context of MRP, it could be
  a simple rule to always first assign the largest process to the least
  loaded machine. The purpose of such approach is most beneficial in
  situations where speed is prioritized over optimality.
\item
  \textbf{Metaheuristic} approaches are higher-level algorithmic
  frameworks designed to guide the underlying heuristics or search
  methods. These provide additional capabilities for escaping local
  optima and are capable of efficiently exploring large solution spaces.
\item
  \textbf{Hyper-heuristic} approaches operate at even higher level than
  metaheuristics, and are capable of automatically generating
  appropriate heuristics or metaheuristics to solve wide range of
  problems, however, also comes at much higher computational cost as
  compared to the other approaches.
\end{itemize}



\subsection{Common Search Methods for tackling Optimization Problems}


\subsubsection{Local Search}

\textbf{Local Search (LS)} method optimizes the objective cost function
through an iterative process of small, incremental improvements. LS
starts from an initial solution and explores the solution space by
applying small modifications, known as \emph{move operators}, to the
current solution. Basic move operators include actions such as process
reassignment (``Shift'') and swapping the assignment of two processes
(``Swap''). Repeated application of these move operators on different
parts of the solution generates a set of possible, ``neighboring''
solutions. Once the solution space is sufficiently explored, the second
component of LS is the \emph{search strategy} to determine the move to
accept; \emph{Steepest descent strategy} generates and evaluates the
entire neighborhood, and moves to the solution with the best improvement
- providing better results at the cost of higher computational cost,
whereas a \emph{Greedy strategy} generates neighboring solutions
one-by-one and immediately accepts the first solution that shows
improvement, even if marginal. The simplest form of a local search
method is known as the \textbf{First-Improvement (Greedy) Hill Climber},
and consists of an LS method using the Greedy search.

\subsubsection{Large Neighbourhood Search}

\textbf{Large Neighborhood Search (LNS)} is a metaheuristic that extends
LS, and consists of two primary components: Destroy phase and Repair
phase. In the first step, Destroy phase, some section of the solution is
``destroyed'', which involves iterative selection of subset of
processes, followed by having them unassigned - for the purpose of
\emph{neighborhood exploration}. The next stage, Repair phase, involves
rebuilding the unassigned processes while optimizing for the cost
function and often includes exact mathematical optimization algorithms,
such as Constraint Programming (CP) or Mixed Integer Programming (MIP),
to find the optimal reassignments for the current process subset. While
these exact solver methods can guarantee optimal solutions, they come
with the trade-off of being computationally expensive due to their
exhaustiveness. \textbf{Variable Neighborhood Search (VNS)} is a
sophisticated variation of LNS, which follows similar processes of
Destroy and Repair phases but with the addition of a diversification
strategy that expands the explored solution space by changing the type
and size of neighborhood during the search process, increasing the
potential for finding a better solution and improved ability to escape
local optima. VNS is also capable of applying intensification strategies
to narrow down the search towards specific neighborhoods through the
adjustment of subproblem selection by using smaller neighborhoods.

\emph{(Not sure if the below are needed - inclusion depends on final
word count)}

\begin{itemize}
\item
  \textbf{Genetic Algorithms}: These algorithms mimic the process of
  natural selection, using techniques such as mutation, crossover, and
  selection to evolve solutions over generations.
\item
  \textbf{Simulated Annealing}: This technique is inspired by the
  annealing process in metallurgy. It explores the solution space by
  allowing occasional uphill moves to escape local optima, with a
  probability that decreases over time.
\item
  \textbf{Tabu Search} method uses a local search procedure and
  maintains a list of previously visited solutions (the ``tabu list'')
  to avoid cycling back to them.
\end{itemize}



\chapter{Related Work}\label{chap:rel_work}

Related Work chapter provides an in-depth analysis of the approaches proposed for tackling the machine reassignment problem, discusses the underlying algorithms and their mechanisms for solution space exploration, search mechanisms, and other techniques implemented alongside the core algorithms. Additionally, commonalities among the solutions are synthesized to provide tabular breakdown categorized into the different components that make the overall solution.



\section{Literature Research Objectives}

The research objectives to be accomplished through literature review are
to gain insights into the different approaches taken by the submitted
solutions for the MRP, investigate whether they were known or novel
techniques, whether the solutions prioritised specific actions over
others, such as preference to initially migrate large processes over
small processes, and whether any sort of randomization played crucial
role in their success. Should common patterns be found, they could
provide insights into the development of future strategies for the
tackling the machine reassignment problem.



\section{Literature Research Method}

The literature research was performed primarily on Google Scholar, NUIG James Hardiman Library, and Scopus for documents containing keywords; "Machine Reassignment Problem", "Google Challenge 2012", "Resource Management", "ROADEF12", and "Resource Optimization".

Abstract sections of the resulting documents from the search queries
were scanned and selected for in-depth analysis depending on their
relevance to the topic. The query was further narrowed down to include
peer-reviewed papers in high-quality journals (Q1 and Q2) that were
published from the year of the publishment of the Google Machine
Reassignment Problem, 2012 and onwards. Furthermore, of the selected
relevant papers, their references and citations were investigated to
provide more informational background and potential implications on the
research that followed.



\section{Proposed algorithms for solving Google Machine Reassignment Problem}

The best performing algorithm in the ROADEF/2012 GMRP contest was presented by team S41 consisting of a variable neighborhood search method including two specialized move operators, a technique for "random" noising of the objective function, and heuristics for process selection. 
\cite{gavranovic2012} 
In second place came in team S38 using LNS as the base for the algorithm and CP solver to compute the optimal reassignments in the context of the selected subset of processes, as well as heuristics for guiding the search and subproblem optimization. 
\cite{mehta2012}
Team J12 placed in third place, and in the first place within the junior category, with an algorithm that switches between two metaheuristics; HC and MIP-LNS. The algorithm initializes with HC and keeps running it until hitting a local optima, then falls back on MIP-LNS, and returns to running HC again following this pattern until the end. 
\cite{jaskowski2016}



\section{Post-contest Contributions}

Due to the overall success of the contest, new algorithms have been
proposed and evaluated using datasets from the GMRP problem for many
years following the ROADEF 2012 contest.



\section{Summary of Findings, Common Patterns, and Insights}

The analysis of the proposed solutions for the MRP revealed a diverse
range of approaches, each with its own strengths and weaknesses.
Heuristic and metaheuristic methods were the most commonly used, likely
due to their ability to provide good solutions within a reasonable
timeframe. Exact methods, while less common, demonstrated the potential
for achieving optimal solutions but were limited by computational
constraints. Common patterns such as initial solution generation, focus
on high-cost processes, iterative improvement, and the use of
randomization were prevalent across many solutions. These patterns
suggest that successful strategies often involve a combination of quick
initial assessments followed by more thorough exploration of the
solution space.

The common characteristics and techniques among the investigated
solutions are: - Use of \emph{heuristics for subproblem selection} with
prioritization based on: (S41) Process Size, (S38) Machine Cost, (J12)
Potential Cost Reduction (Load \& Balance), (J12) Potential Objective
Cost Reduction - Incorporation of \emph{randomness}: (S41) Noising of
the objective function, (S38) Random machines in the subproblem
selection, (J12)



\chapter{Methodology} \label{chap:methodology} 



\section{Dataset}

The datasets containing the initial process-machine assignments, as well as other variables described in the problem definition, are provided on the official ROADEF/EURO 2012 challenge page.

Only one dataset, A, is used for the analysis of reassignments as it contains a smaller set of instances, allowing the algorithms to complete more reassignments as compared to the more difficult datasets, B and X. Given the underlying logic for reassignments remains the same, it is assumed the reassignment behavior will also remain the stay, whilst at
the same time a larger dataset that will be generated from tacking the reassignments is preferable.

Overview of datasets and the specific variables of the two relevant dataset A instances are:

\begin{table}[h]
    \centering
    \begin{tabular}{l|c|c|c|c}
        Variable & Dataset A & Datasets B and X & Instance a1\_1 & Instance a1\_4 \\ 
        \hline
        Processes      & \([100, 1000]\)   & \([5000, 50000]\) & 4  & 50   \\ 
        Machines       & \([4, 100]\)      & \([100, 5000]\)   & 50 & 1000 \\ 
        Resources      & \([2, 12]\)       & \([3, 12]\)       & 2  & 3    \\
        Transient      & \([0, 4]\)        & \([0, 4]\)        & 0  & 1    \\ 
        Services       & \([10, 125]\)     & \([500, 1000]\)   & 10 & 100  \\ 
        Neighbourhoods & \([1, 50]\)       & \([5, 5]\)        & 1  & 50   \\ 
        Locations      & \([1, 50]\)       & \([10, 100]\)     & 4  & 50   \\ 
        Dependencies   & \([0, 10]\)       & \([30, 70]\)      & 10 & 10   \\ 
        Balance        & \([0, 1]\)        & \([0, 1]\)        & 1  & 1    \\ 
    \end{tabular}
    \caption{Example dataset characteristics}
\end{table}



\section{Algorithms Selection for Analysis of Reassignments}

For an in-depth analysis into the algorithms and their reassignments
throughout the optimization process, the three best performing,
state-of-the-art algorithms are selected, namely: - VNS proposed by team
S41, - CP-LNS proposed by team S38, - HC-(MIP-LNS)-HC proposed by team
J12

VNS-based approach incorporates two specialized move operators, along
side the two basic operators, used for the exploration of the
neighborhood; - \texttt{Chain} involves triggering a sequence of
interdependent reassignments, such as cyclic exchanges, with the purpose
of further expanding the exploration space - \texttt{BPR} (Big Process
Reassignment) specifically targets the reassignment of large processes
by removing multiple of processes at once, as these are generally the
hardest to reallocate, and improves the efficiency of the other three
operators (Shift, Swap, Chain) that often struggle with reassignment of
the larger processes. BPR is also attributed the highest potential of
contribution towards the improvement of objective cost. Additionally,
the search is guided by heuristic to sort processes by size and
prioritize the reassignment of large processes at first, as it becomes
harder to do so in later stages, and explores the neighborhood using
move operators in the following order; BPR, Shift, Swap, Chain. Another
technique incorporated into the algorithm aims to expand the
diversification strategies and consists of ``noising'' the objective
function, which temporarily increases the weight cost of a constrained
resource, and then optimizing for the noised objective function, forcing
the search out of local optima. The use of specialized BPR operator as
well as process size heuristic indicates this approach could be a good
candidate to extract meaningful information based on reassignments,

In \cite{mehta2012}, both MIP and CP solvers were evaluated with an LNS-based
approach. For MIP, performance gains were observed in the selection of
all process from small number of machines over selection of few
processes from a large number of machines. However, CP has shown better
performance on larger instances and was selected as the solver for the
final submission. Additionally, significant gains that pushed the base
CP-LNS performance are attributed to multiple domain-specific heuristics
at different layers of the algorithm. Specifically, cost-driven
heuristic is used to guide the search by prioritizing machines with the
highest current cost to be used for the subproblems, mixed with some
machines selected at random - with in this approach, only a subset of
processes of the selected machines is considered for reassignment to
reduce probability of situations where a process is assigned back to the
same machine from which it was just unassigned. The variable domains of
CP are also updated using heuristics to improve performance of
subproblem optimization by reducing the search space; (1) \emph{variable
ordering heuristic} are based on aggregation of potential improvement of
objective cost, the total weight of process requirements, as well as the
available machines for the process, (2) \emph{value ordering heuristics}
define the selection of suitable machine for a process based on minimum
cost, and (3) \emph{filtering rules} to narrow down the solution space
by removing machines from consideration that would produce infeasible
solutions. Based on the heuristics in this approach, it is difficult to
make assumptions on any expectations for potential patterns in
reassignments as the cost-driven heuristic does not target processes of
specific size, but instead machines of specific cost - where even
knowing a machine has high cost, an indication of containing hosting
something heavy, could mean it is either many small processes or few
large processes.

Even though MIP solver was observed to provide suboptimal performance on
larger instances based on investigations of the previous approach, the
HC-LNSHC algorithm is capable of delivering solid performance with the
MIP solver due to its architectural design consisting of two
metaheuristics working together; (1) Local Search with Greedy search,
i.e.~First-improvement Hill Climber (FI-HC), and (2) LNS using MIP for
subproblem optimization. Rather than relying on the computationally
expensive exact solvers on each iteration, the algorithm instead
initializes with FI-HC to rapidly iterate towards any improvement in the
objective cost as fast as possible. The iteration speed is further
optimized with an ``Optimistic Improvement'' heuristic for subproblem
selection based on estimation of the potential for cost reduction
derived from load and move costs, and then prioritizing processes with
the highest potential - rather than having to do the more expensive
computation of entire solution cost. In addition, the FI-HC also
maintains a tabu list of the previously explored machines to avoid
cycling. The FI-HC is left running until eventually reaching getting
stuck in a local optima, from which if it is unable to find further
improvements, the algorithm switches to MIP-LNS. The MIP-LNS also uses
heuristics for subproblem selection to prioritize subsets based on
potential to reduce the objective cost, while the subproblem size is
adaptively adjusted by starting off small, and iteratively increasing
the subproblem size if it is optimally solved until reaching a point
where the optimization fails, at which point the previously solved
subproblem is accepted as the best. After the MIP-LNS completed, the
algorithm falls back to the FI-HC again and repeats the process. Based
on the analysis of the approach, the performance can be attributed to
the overall design of the algorithm and the effectiveness of heuristics
incorporated into both the FI-HC and MIP-LNS. It is possible the use of
MIP on larger instances is also resulting in similar performance hits
when ran that were observed in previous approach, but the FI-HC could be
minimizing the impact through its efficient and fast iterations.

Overall, all three algorithms incorporate some level of domain-specific
heuristics that provide significant performance gains well above
original capabilities.



\section{Metrics for the analysis of reassignments}

No datasets of reassignments are generated by the existing algorithms
and therefore requires the design and implementation of a tracking
system into the algorithms to generate the dataset on reassignments. The
metrics tracked for each reassignment should provide sufficient
information for the analysis and include:

\begin{itemize}
\item
  \textbf{Process ID}: Serves as the primary ID to identify the process
  that is being reassigned at the current step
\item
  \textbf{Process Size}: Provides information on the aggregation of
  process resource requirements to capture insights into what types of
  processes are being reassigned, how frequently, and potential
  patterns.
\item
  \textbf{Solution ID}: Allows to derive further information on how the
  solution evolves over time such as the number of reassignments
  performed in the process of moving from the current solution to the
  next solution accepted by the search mechanism.
\item
  \textbf{Total Objective Cost}: Associates reassignments and their
  potential quality through the impact on the objective function and
  should allow to track the evolution of objective cost over
  reassignments.
\item
  \textbf{Improvement Cost}: This is the delta between the objective
  cost between two adjacent reassignments. It is primary tracked for
  convenience to provide insights into the direction and size of
  improvement in objective cost for each reassignment, to tackle
  questions such as whether the reassignment of processes of specific
  size (large vs.~small) has greater potential.
\end{itemize}



\section{Tracking System Design}

The goal of the tracking system is to provide the ability of capturing
reassignments and achieved by modifying the original codebases of each
algorithm and have them write the reassignment metrics to a file each
time a reassignment is performed. This approach is not optimal in terms
of algorithm performance, as writing to file and especially this
frequently for every assignment, is an expensive operation and will slow
down the algorithms. However, the approach is still suitable for the
analysis as the performance impact does not affect the types of
reassignments, just their number - even with less recorded
reassignments, it is more than enough.

Steps taken in the planning and design of the tracking system:

\begin{enumerate}
\def\labelenumi{\arabic{enumi}.}
\setlength\itemsep{0em}
\setlength\parskip{0pt}
\setlength\parsep{0pt}
\item Clone repositories containing the algorithms,
\item Run the algorithms before making any changes to ensure they are working as expected,
\item Investigate structure of the project to understand the implementation of the algorithms,
\item Design an approach for the implementation of the tracking system,
\item Modify the original codebases to capture reassignments,
\item Run the algorithms using datasets of initial assignments and capture the reassignments,
\item Finally, use the newly generated reassignments dataset for the remaining analysis.
\end{enumerate}

Source code for the original algorithms was found at:

\begin{itemize} \setlength\itemsep{0em} \setlength\parskip{0pt} \setlength\parsep{0pt}
    \item S41 (VNS): https://github.com/harisgavranovic/roadef-challenge2012-S41
    \item S38 (CP-LNS): https://sourceforge.net/projects/machinereassign
    \item J12 (HC-LNSHC): https://bitbucket.org/wjaskowski/roadef-challange-2012-public
\end{itemize}

The modified codebases to track reassignments are available under one
repository:

- Link: https://github.com/yourusername/reassignment-tracking



\section{Tracking System Implementation Challenges and Limitations}

The first codebase investigated was of the HC-LNS-HC algorithm,
developed in Java. The main difficulty, and a common theme across all
three codebases, is the lack of information on the commands or arguments
required for running the projects. The LNS is using MIP solver, which at
the time of original implementation was CPLEX Solver v12.5 - however,
only the latest version is available for download, which resulted in the
requirement for updating the dependency to CPLEX v22.1.1, and along with
it multiple additional changes to the codebase as the change new CPLEX
version is not backwards compatible. Major challenges were encountered
in the following codebases of CP-LNS and VNS algorithms, primarily due
to all algorithms being implemented in different programming languages
(Java, C, C++), preventing the ability of reusing code for tracking
assignments and instead requiring individual implementations each time.
Additionally, the CP-LNS algorithm is implemented in a way that makes
retrieving much of the reassignment metrics not possible - such as the
process size which in the end required post processing steps on the
captured reassignments to populate the column using dataset generated by
one of the other algorithms, since all algorithms were ran using the
same instance. Secondly, the approaches are lacking the concept of
``Solution ID'' and this metric was not feasible to be tracked, as well
as it performing both the explored and accepted reassignments in single
\texttt{for} loop, making the tracking of actual reassignments very
difficult. Unfortunately the inclusion of explored reassignments during
neighborhood exploration is not feasible as it generates a reassignments
file of over 16GB and crashes the Python kernel, and required
significant efforts to complete the implementation while ensuring
exploration moves are not tracked. The tracking implementation in the
CP-LNS codebase is questionable to some extent as it relies on flipping
a flag back and forth in an attempt to filter out exploration moves and
only capture true reassignments.



\chapter{Experimental Setup} \label{chap:experiments}


\section{Hardware specifications and system dependencies}

The algorithms were executed on the same hardware of following specifications:

\begin{itemize} \setlength\itemsep{0em} \setlength\parskip{0pt} \setlength\parsep{0pt}
    \item 13th Gen Intel Core i9-13950 HX with 24 cores (32 threads) and max turbo frequency up to 5.50 GHz.
    \item 192 GB DDR4 RAM @ 3600 MT/s
\end{itemize}

In addition, the following system dependencies are required to support
the execution of investigated algorithms:

\begin{itemize} \setlength\itemsep{0em} \setlength\parskip{0pt} \setlength\parsep{0pt}
    \item  Ubuntu 22.1 - but should also work without changes on other UNIX-based systems and will require changes to run on Windows OS as some build scripts are written in bash and would need to be converted to PowerShell.
    \item JDK with Java 8 support
    \item Apache Ant
    \item IBM ILOG CPLEX Studio v22.1.1 (specifically for the MIP solver)
    \item GCC C++ compiler with C++11 support
    \item Make
    \item Python 3.x
\end{itemize}



\section{Streamlined pipeline to run the algorithms}

VS Code was used in both the implementation and analysis stages of the project as it provides support for all necessary tools and technologies used through the investigation.

During the implementation stage, a more streamlined pipeline was developed for testing and running the algorithms with the use of "Run Tasks" feature provided in VS Code. The feature allows to define set of configurations in a \texttt{tasks.json} file such as the commands to build and run the algorithms, as well as to pass in any additional arguments. These tasks resulted in significant improvement in convenience throughout the investigation as not the algorithms are all implemented in different programming languages, and hence require different commands to run, but also multiple arguments need to be passed into the program including; path to initial assignments dataset, path to file containing data on the problem variables (constraints), output file, and other project specific arguments. With the current pipeline, any algorithm can be executed by simply opening the Command Palette in VS Code Run Tasks \textgreater{} and selecting the specific task to execute, e.g. "HC-LNSHC: Run Full Duration Analysis (a1\_4, 300s)" or "HC-LNSHC: Test Run (a1\_1, 10s)" to either run the algorithm for the full 5 minute duration or a quick test run for 10 seconds on a different, easier dataset that allow for faster development iteration cycles throughout the implementation and testing of the changes for tracking reassignments.





\chapter{Results} \label{chap:results}

This chapter presents the results of the experimental runs conducted to evaluate the selected algorithms, with the primary focus into the analysis of reassignments made by the algorithms during their execution. Once the plots are introduced, they are discussed in detail to highlight the key findings and observations.

\begin{table}[h]
    \centering
    \begin{tabular}{l|c|c|c}
        Variable & VNS & CP-LNS & HC-LNSHC \\ 
        \hline
        #Reassignments      & \([100, 1000]\)   & \([5000, 50000]\) & 4  & 50   \\ 
        Unique Processes       & \([4, 100]\)      & \([100, 5000]\)   & 50 & 1000 \\ 
        #Neighbourhoods      & \([2, 12]\)       & \([3, 12]\)       & 2  & 3    \\
    \end{tabular}
    \caption{Summarised characteristics of reassignments of the three algorithms.}
\end{table}

\begin{figure}[htbp]
  \centering
  \includegraphics[width=135mm]{/home/pbiel/repos/mrp/analysis/images/hclns_reassignments.png}
  \caption{HC-LNSHC Algorithm Process Sizes of each Reassignment}
  \label{fig:hclns_reassignments}
\end{figure}

Presents the reassignments made by the HC-LNSHC algorithm during its execution on the dataset instances a1\_1 and a1\_4. The plots illustrate the number of reassignments over time, providing insights into the algorithm's behavior and performance.

\begin{figure}[htbp]
  \centering
  \includegraphics[width=135mm]{/home/pbiel/repos/mrp/analysis/images/vns_reassignments.png}
  \caption{VNS Algorithm Process Sizes of each Reassignment}
  \label{fig:vns_reassignments}
\end{figure}

Presents the reassignments made by the VNS algorithm during its execution on the dataset instances a1\_1 and a1\_4. The plots illustrate the number of reassignments over time, providing insights into the algorithm's behavior and performance.

\begin{figure}[htbp]
  \centering
  \includegraphics[width=135mm]{/home/pbiel/repos/mrp/analysis/images/cplns_reassignments.png}
  \caption{CP-LNS Algorithm Process Sizes of each Reassignment}
  \label{fig:cplns_reassignments}
\end{figure}

Presents the reassignments made by the CP-LNS algorithm during its execution on the dataset instances a1\_1 and a1\_4. The plots illustrate the number of reassignments over time, providing insights into the algorithm's behavior and performance.



\chapter{Conclusion} \label{chap:conclusion}

The main challenge with the available solutions stems from the lack of
ability to evaluate based purely on their conceptual approaches and
improved evaluation approach of the solutions would be the
standardization of technologies and tools, such to provide even grounds
for evaluation. This is due to the fact that an algorithm programmed in
Java comes with the overhead of the JVM as well as of the garbage
collector - whereas, the other two approaches programming in C and C++
do not.







%%%%%%%%%%%%%%%%%%%%%%%%%%%%%%%%%%%%%%%%%%%%%%%%%%%%%%%%%%%%%%%%%%%%%%%%%%%%%%%%
%\bibliographystyle{plainnat}                  % to give author-year style
\bibliographystyle{IEEEtranN} 
\renewcommand{\bibname}{References}           % change default name Bibliography to References
\bibliography{references}                     % BibTeX References file, references.bib
\addcontentsline{toc}{chapter}{References}    % add References to TOC


%%% uncomment if Appendices needed
%\appendix
%\chapter{Appendix-A-Title} 
%\label{chap:appendix_a}

%\chapter{Appendix-B-Title} 
%\label{chap:appendix_b}
	
\end{document}
